\section{本章小结}
我们提出了一种新颖的模态综合法（Component Mode Synthesis, CMS），解决了传统子结构技术在大规模特征值问题中收敛缓慢的问题。通过识别固定界面子结构模态中的相位缺失，我们引入了一种多重划分策略，构建了相位完备的子空间。此外，我们开发了一种无 Schur 补和无特征求解的形式，显著降低了界面模态的计算开销和内存占用。对比实验和可扩展性测试表明，我们的方法相比之前的 CMS 变体，精度提高了多达三个数量级，同时保留了快速局部更新和进程间通信量小等关键优势。

目前我们的工作存在一些局限性。首先，该方法主要关注最低频的特征模态；将其扩展到目标位移附近的内部特征值问题仍是一个开放性问题。其次，当前的划分策略——依赖于简单的交错划分或启发式的 METIS 调整——可能并非最优。探索由误差估计器驱动的自适应划分方案可能进一步提升性能。

展望未来，研究如何将我们的多重划分概念与 AMLS \cite{Bennighof2004AMLS} 等多层级框架相结合是一个有前景的方向。这种混合方法有望带来乘积级的效率提升，前提是能够解决跨层级相位自适应性的挑战。